\section{Synchronisation and Dependences}
\label{sec:sync}

The previous sections presented DSTM as a memory system. This one presents it
as a \emph{schedule}, in the vocabulary of dependence analysis and scheduling.
Nothing new is assumed: everything follows from
\crefrange{res:monotone}{res:terminate}. The payoff is that the results become
graph-theoretic and the design's structure becomes visible.

% ===========================================================================
\subsection{Bernstein conditions over versioned cells}
\label{sec:sync-bernstein}

Bernstein's conditions \cite{bernstein1966} state when two program fragments
may be reordered: when neither writes what the other touches. With footprints
as in \cref{def:footprint}:

\begin{definition}[Independence]\label{def:independence-ops}
Operations $o_1, o_2$ are \emph{independent}, written
$o_1 \disjoint o_2$, when
\[
  \underbrace{\fpw{o_1} \cap \fpr{o_2} = \emptyset}_{\text{no flow (RAW)}}
  \quad\wedge\quad
  \underbrace{\fpr{o_1} \cap \fpw{o_2} = \emptyset}_{\text{no anti (WAR)}}
  \quad\wedge\quad
  \underbrace{\fpw{o_1} \cap \fpw{o_2} = \emptyset}_{\text{no output (WAW)}} .
\]
\end{definition}

The decisive structural fact is that a footprint carries an effective time.

\begin{theorem}[Versioning is renaming]\label{thm:renaming}
If $\opeff{o_1} \ne \opeff{o_2}$ then $o_1 \disjoint o_2$, whatever their
\glspl{rank}, addresses and lengths.
\end{theorem}

\begin{proof}
By \cref{def:footprint}, $\fp{o_i} \subseteq \{\optgt{o_i}\} \times
\{\opeff{o_i}\} \times \Addr$. \Glspl{cell} carry their effective time as a
coordinate, so footprints at distinct \glspl{effective-time} are disjoint, and all
three intersections of \cref{def:independence-ops} are empty.
\end{proof}

\Cref{thm:renaming} is short but is why DSTM needs neither locks nor rollback.
In a conventional \gls{darray}, overwriting a location at tick $t+1$ creates
an output dependence on the write at tick $t$ and an anti-dependence on every
in-flight read of tick $t$; enforcing those is what a per-tick \gls{barrier}
does. DSTM removes them by \emph{renaming}: the tick index is part of the
location's name, so the two writes are not to the same location. This is the
manoeuvre of static single assignment \cite{cytron1991ssa} and array expansion
\cite{feautrier1991}, applied along a simulation-time axis.

What survives is only what happens \emph{within} one effective time.

\begin{theorem}[Residual dependences]\label{thm:residual-dep}
Let $\Hist$ be legal. Then for all $o_1 \ne o_2$ in $\Hist$, either
$o_1 \disjoint o_2$, or $o_1$ and $o_2$ are a put and a get with
$\optgt{o_1} = \optgt{o_2}$, $\opeff{o_1} = \opeff{o_2}$ and overlapping
addresses --- a flow dependence from the put to the get.
\end{theorem}

\begin{proof}
By \cref{thm:renaming} assume $\opeff{o_1} = \opeff{o_2} = \eff$ and
$\optgt{o_1} = \optgt{o_2}$, otherwise independence is immediate. Three cases.

\emph{Both \glspl{put}.} Then $\fpw{o_1} \cap \fpw{o_2} = \emptyset$ by
\cref{res:write-disjoint}, so there is no output dependence.

\emph{Both \glspl{get}.} $\fpw{o_i} = \emptyset$ for both, so all three intersections
are empty.

\emph{One put $w$, one get $g$.} The only possibly non-empty intersection is
$\fpw{w} \cap \fpr{g}$, a flow dependence. The anti-dependence
$\fpr{w} \cap \fpw{g}$ is empty because $\fpr{w} = \fpw{g} = \emptyset$ by
definition.
\end{proof}

\begin{corollary}[Normal form]\label{cor:dsa}
Under legality, a DSTM execution is in dynamic single assignment form: every
cell is written at most once, and the only dependences are flow dependences
from that unique write to its readers. Anti- and output dependences do not
occur at any granularity.
\end{corollary}

Two things eliminate them. Across effective times \cref{thm:renaming} does so
\emph{structurally} --- no restriction is involved and no program can
reintroduce them. Within one effective time \cref{res:write-disjoint}
eliminates output dependences and \cref{res:phase} anti-dependences, both
program restrictions a bad program can violate. DSTM's determinism is thus
partly free and partly earned.

% ===========================================================================
\subsection{The dependence graph}
\label{sec:sync-dag}

\begin{definition}[Dependence graph]\label{def:depgraph}
The \emph{dependence graph} $\depgraph(\Hist)$ has
\begin{itemize}[leftmargin=1.5em,itemsep=1pt]
  \item a node $\task{p}{t}$ for each task (\cref{def:task}) of $\Hist$, with
    node weight $\work{p}{t}$;
  \item a \emph{sequential} edge $\task{p}{t} \to \task{p}{t'}$ whenever $t <
    t'$ are consecutive \glspl{origin-time} used by $p$;
  \item a \emph{flow} edge $\task{p}{t} \to \task{r}{t'}$ whenever a put
    issued by $p$ at origin $t$ and a get issued by $r$ at origin $t'$ stand
    in the flow dependence of \cref{thm:residual-dep}.
\end{itemize}
\end{definition}

\begin{lemma}[Edge span]\label{lem:edge-span}
Let $\task{p}{t} \to \task{r}{t'}$ be a flow edge arising from a put with
\gls{delta} $\delta$ and a get with delta $\delta'$, both at effective time $\eff$
on owner $q$. Then
\[
  t' - t \;=\; \delta + |\delta'| ,
\]
and $t' - t \ge 1$ whenever $p \ne r$.
\end{lemma}

\begin{proof}
$t = \eff - \delta$ and $t' = \eff - \delta' = \eff + |\delta'|$, giving the
identity. For the bound, suppose $p \ne r$. If $p \ne q$ then
\cref{res:lookahead} and (M2) give $\delta \ge \dist{p}{q} \ge 1$. If $p = q$
then $r \ne q$, so the get is remote and the rule \textsc{Get-Remote} requires
$\delta' \le -1$, giving $|\delta'| \ge 1$. Either way $t' - t \ge 1$.
\end{proof}

\begin{theorem}[Schedulability]\label{thm:acyclic}
For a legal history, $\depgraph(\Hist)$ is acyclic.
\end{theorem}

\begin{proof}
Order nodes by $\task{p}{t} \prec \task{r}{t'}$ iff $t < t'$, or $t = t'$ and
the operations of $\task{p}{t}$ precede those of $\task{r}{t'}$ in the program
order of the common process. Sequential edges increase $t$ by construction. Flow
edges with $p \ne r$ increase $t$ strictly by \cref{lem:edge-span}. Flow edges
with $p = r$ have $t' - t = \delta + |\delta'| \ge 0$, and when $t' = t$ both
$\delta$ and $\delta'$ are zero, so the put and the get are local, at the same
effective time, on the same process --- and \cref{res:phase} orders the put
before the get in \gls{program-order}. Every edge therefore strictly increases
$\prec$, which is a strict partial order, so no cycle exists.
\end{proof}

\Cref{thm:acyclic} is the graph-theoretic form of \cref{thm:determinism}: the
induction used there is a topological order of $\depgraph$. It is also that of
\cref{thm:deadlock-freedom}, since a deadlock is a cycle in the wait-for
graph, itself a sub-graph of $\depgraph$ reversed.

\begin{remark}[Recovering the Chandy--Misra condition]\label{rem:chandy-misra}
\Cref{lem:edge-span} is the only place positivity is used, to show every
cross-rank edge spans at least one tick. Conservative PDES states the same
requirement as ``no cycle of zero lookahead'' \cite{chandy1979}; here it is a
property of $\cgraph$, holding by construction rather than by analysis, which
is what lets DSTM dispense with null messages.
\end{remark}

% ===========================================================================
\subsection{Push and pull are the same schedule}
\label{sec:sync-duality}

\begin{figure}[htbp]
\centering
\resizebox{\textwidth}{!}{%
\begin{tikzpicture}[
  x=20mm, y=-13mm,
  every node/.style={font=\scriptsize},
  tk/.style={draw,rounded corners=1.5pt,minimum width=13mm,minimum height=6.5mm,
             fill=resbg,draw=dstmblue,inner sep=1pt},
  vv/.style={draw,rounded corners=1.5pt,minimum width=13mm,minimum height=6.5mm,
             fill=white,draw=dstmgrey,dashed,inner sep=1pt},
  >={Stealth[round]},
]
  % ---- pull ----
  \node[font=\small\bfseries,text=dstmblue,anchor=west] at (-0.9,-0.9) {Pull};
  \node[anchor=east,text=dstmgrey] at (-0.35,0) {$p$};
  \node[anchor=east,text=dstmgrey] at (-0.35,1) {$r$};
  \node[tk] (pa) at (0,0) {$t$};
  \node[vv] (va) at (1,0) {$\ver{p}{t}$};
  \node[tk] (ra) at (2,1) {$t{+}D$};
  \draw[->,thick,dstmblue] (pa) -- node[above,inner sep=1.5pt]{$\opput\ \delta{=}0$} (va);
  \draw[->,thick,dashed,dstmred] (va) -- node[above right,inner sep=1.5pt]{$\opget\ \delta{=}{-}D$} (ra);

  % ---- push ----
  \node[font=\small\bfseries,text=dstmblue,anchor=west] at (-0.9,2.1) {Push};
  \node[anchor=east,text=dstmgrey] at (-0.35,3) {$p$};
  \node[anchor=east,text=dstmgrey] at (-0.35,4) {$r$};
  \node[tk] (pb) at (0,3) {$t$};
  \node[vv] (vb) at (1,4) {$\ver{r}{t{+}D}$};
  \node[tk] (rb) at (2,4) {$t{+}D$};
  \draw[->,thick,dstmblue] (pb) -- node[above right,inner sep=1.5pt]{$\opput\ \delta{=}{+}D$} (vb);
  \draw[->,thick,dashed,dstmred] (vb) -- node[above,inner sep=1.5pt]{$\opget\ \delta{=}0$} (rb);

  % ---- common conclusion ----
  \draw[decorate,decoration={brace,amplitude=4pt,mirror},dstmamber,thick]
       (2.42,-0.35) -- (2.42,4.35)
       node[midway,right=6pt,text=dstmamber,align=left,font=\scriptsize]
       {consumer tick is\\$t + D$ either way};
\end{tikzpicture}}
\caption{The two idioms of \cref{sec:exec-idioms}. Solid edges are writes,
dashed edges reads, dashed boxes versions. The version they route through
differs --- $p$'s own region versus $r$'s --- but the tick at which the value
reaches the consumer is $t + \dist{p}{r}$ in both cases.}
\label{fig:push-pull}
\end{figure}

\begin{theorem}[Pipeline depth]\label{thm:pipeline-depth}
Let $p$ produce a value at origin time $t$ and let $r$ consume it, with
$q \in \{p, r\}$ the owner of the mediating version, $\dist{p}{r} = D$.
\begin{enumerate}[label=(\roman*),leftmargin=1.4cm,itemsep=1pt]
  \item In the \glslink{pull-model}{pull idiom} ($q = p$, $\delta = 0$, $\delta' = -D$) the consumer
    tick is $t' = t + D$.
  \item In the \glslink{push-model}{push idiom} ($q = r$, $\delta = +D$, $\delta' = 0$) the consumer
    tick is $t' = t + D$.
\end{enumerate}
Moreover $t' = t + D$ is optimal: no legal choice of $\delta, \delta'$ gives
$t' < t + D$.
\end{theorem}

\begin{proof}
Both cases are \cref{lem:edge-span}: $t' - t = \delta + |\delta'|$, which is
$0 + D$ and $D + 0$ respectively. For optimality, if $q = p$ then
$\delta \ge 0$ and, since $r \ne q$, \cref{res:staleness} bounds
$|\delta'| \le \dist{r}{p} = D$ while \textsc{Get-Remote} forces
$|\delta'| \ge 1$; the minimum of $\delta + |\delta'|$ subject to the reader
actually reaching $p$'s version is attained at $\delta = 0$, and the consumer
tick can be no earlier than the tick at which the reader may legally look back
to $t$, namely $t + |\delta'|$ with $|\delta'| \le D$. Symmetrically if
$q = r$ then $\delta \ge \dist{p}{r} = D$ and $\delta' = 0$, so
$t' - t \ge D$. In both cases $t' - t \ge \min(D, D) = D$ once the constraint
that the value must actually cross from $p$ to $r$ is imposed.
\end{proof}

\begin{remark}[Push or pull is a cost decision]
The corollary is easy to miss: the choice between push and pull is never
semantic and never affects how far processes may drift. It is a
communication-cost decision. Push moves one message per producer--consumer
pair per tick, pull two; push is the better default and is what both benchmark
applications in \cite{dstm2026} use. Pull earns its cost when the consumer
reads only occasionally, moving data on demand.
\end{remark}

% ===========================================================================
\subsection{Slack, drift and the critical path}
\label{sec:sync-slack}

Since $\depgraph$ is a weighted DAG, the standard scheduling quantities are
available.

\begin{definition}[Earliest start, slack]\label{def:slack}
For a node $u$ of $\depgraph$, let $\mathit{est}(u)$ be the length of the
longest weighted path ending at $u$ (excluding $u$'s own weight), and
$\mathit{lst}(u)$ the latest start compatible with the overall makespan. The
\emph{slack} is $\slack{}{u} \defeq \mathit{lst}(u) - \mathit{est}(u)$.
\end{definition}

Drift (\cref{def:drift}) realises slack: a process runs ahead exactly as far
as its tasks have slack, and stops --- in a blocking get --- on reaching a
task with none. A barrier-synchronised execution zeroes all slack at every
tick boundary, which is why it is slower and why its cost is set by the
slowest process.

\begin{theorem}[Makespan]\label{thm:makespan}
Let $\cpath(\Hist)$ be the maximum, over directed paths $\pi$ in
$\depgraph(\Hist)$, of $\sum_{u \in \pi} \work{}{u}$. Then, under
\cref{def:admissible} and with zero communication cost,
\[
  \makespan_{\mathrm{DSTM}} \;=\; \cpath(\Hist)
  \qquad\text{and}\qquad
  \makespan_{\mathrm{BSP}} \;=\; \sum_{t} \max_{p} \work{p}{t},
\]
and $\makespan_{\mathrm{DSTM}} \le \makespan_{\mathrm{BSP}}$.
\end{theorem}

\begin{proof}
The first equality is the standard characterisation of a greedy schedule on a
DAG with unbounded processors, available because \cref{thm:acyclic} gives
acyclicity and each process executes its own tasks in order. For the inequality,
observe from the proof of \cref{thm:acyclic} that every edge strictly
increases the origin time or, at equal origin time, the program-order
position; hence a directed path visits at most one node per origin time $t$.
Therefore
\[
  \sum_{u \in \pi} \work{}{u}
  \;\le\; \sum_{t} \max_{p} \work{p}{t}
  \;=\; \makespan_{\mathrm{BSP}} . \qedhere
\]
\end{proof}

\begin{remark}[Where the measured speedup comes from]
The gap between the two expressions in \cref{thm:makespan} is what the
evaluation in \cite{dstm2026} reports. It is large when per-tick load is
imbalanced but not persistently so: when the expensive process differs from
tick to tick, $\sum_t \max_p \work{p}{t}$ sums a sequence of maxima while the
critical path weaves between processes. It vanishes when one process is
uniformly slowest, the critical path then running down that row.
\end{remark}

% ===========================================================================
\subsection{Memory as a function of drift}
\label{sec:sync-space}

Versioning trades memory for synchronisation. The trade is bounded, which is
why the scheme is practical.

\begin{theorem}[Space bound]\label{thm:space-bound}
In any reachable configuration, the number of live versions at process $\proc{p}$
satisfies
\[
  \bigl|\{\eff : \Vst{p}{\eff} \ne \Absent\}\bigr|
  \;\le\; H_p - L_p + 1,
\]
where $H_p$ is the largest effective time targeted by any put issued so far
towards $p$, and
\[
  L_p \;\defeq\; \min_{q \in \nbr{p}} \bigl(\kn{p}{q} - \dist{p}{q}\bigr),
  \qquad L_p \defeq \clk{p} \text{ when } \nbr{p} = \emptyset .
\]
\end{theorem}

\begin{proof}
A version exists only if created by \textsc{Apply}, so no live version exceeds
$H_p$. For the lower end, any version with $\eff < L_p$ satisfies
$\kn{p}{q} - \dist{p}{q} > \eff$ for every $q \in \nbr{p}$; together with
$\clk{p} \ge \eff$ (which holds since $L_p \le \clk{p}$ by
\cref{inv:l1-clock}) and the absence of pending gets --- guaranteed by
\cref{thm:discard-safety}, since such a get would be illegal --- the premises
of \textsc{Retire} hold, and \textsc{Evict} follows by \ref{f:progress}.
\end{proof}

\begin{corollary}[Memory is linear in drift]\label{cor:space-drift}
Suppose the execution has bounded drift $\Theta$ (\cref{def:bounded-drift}),
the horizon bound of \cref{res:horizon} is $\Theta_{\max}$, and heartbeats are
delivered promptly enough that $\kn{p}{q} \ge \clk{q} - \lambda$ for a
constant $\lambda$. Then
\[
  \bigl|\{\eff : \Vst{p}{\eff} \ne \Absent\}\bigr|
  \;\le\; 2\Theta + \Theta_{\max} + \diam_p + \lambda + 1 .
\]
In particular the live-version count is independent of the length of the run.
\end{corollary}

\begin{proof}
$H_p \le \max_{q \in \nbr{p} \cup \{p\}} \clk{q} + \Theta_{\max}
 \le \clk{p} + \Theta + \Theta_{\max}$ by bounded drift and
\cref{res:horizon}. And
$L_p \ge \min_q (\clk{q} - \lambda) - \diam_p
 \ge \clk{p} - \Theta - \lambda - \diam_p$. Subtract.
\end{proof}

\Cref{cor:space-drift} is the formal content of the claim in \cite{dstm2026}
that discarding \glspl{version} ``is a memory management optimization $\dots$
without changing the semantics of the program''. Both halves are theorems:
semantics-preservation is \cref{thm:discard-safety}, the bound is
\cref{cor:space-drift}. Each term says something operational --- $2\Theta$
that memory is the price of \gls{drift}, $\Theta_{\max}$ that writing far into
the future costs, $\diam_p$ that distant neighbours cost, $\lambda$ that a
lazy \gls{heartbeat} implementation costs memory rather than correctness.

% ===========================================================================
\subsection{Summary}
\label{sec:sync-summary}

\begin{table}[htbp]
\centering
\small
\begin{tabularx}{\textwidth}{@{}>{\raggedright\arraybackslash}p{4.1cm}
                                >{\raggedright\arraybackslash}X l@{}}
\toprule
\textbf{DSTM mechanism} & \textbf{Dependence-theoretic reading}
 & \textbf{Result} \\
\midrule
Effective time in the cell identity
 & renaming; kills WAR and WAW across ticks
 & \cref{thm:renaming} \\
\Cref{res:write-disjoint}
 & dynamic single assignment within a tick
 & \cref{cor:dsa} \\
\Cref{res:phase}
 & write-phase before read-phase per version; kills WAR within a tick
 & \cref{thm:residual-dep} \\
Delta of an operation
 & weight of a producer--consumer edge
 & \cref{lem:edge-span} \\
(M2), $\dist{p}{q} \ge 1$
 & every cycle of $\cgraph$ has positive weight
 & \cref{thm:acyclic} \\
Distance matrix
 & declared lookahead; pipeline depth
 & \cref{thm:pipeline-depth} \\
Process drift
 & slack in the schedule
 & \cref{def:slack} \\
Absence of a per-tick barrier
 & critical path instead of sum-of-maxima
 & \cref{thm:makespan} \\
Version retirement
 & liveness analysis; storage reclamation
 & \cref{thm:space-bound} \\
Three-round ring barrier
 & termination detection
 & \cref{thm:barrier-correct} \\
\bottomrule
\end{tabularx}
\caption{The dictionary between DSTM's mechanisms and their
dependence-theoretic content.}
\label{tab:sync-dictionary}
\end{table}
